\documentclass[UTF8,12pt,twoside]{ctexart}

\usepackage{myTemp}
\lhead{样例模板}

\begin{document}
%%%%%%%%%%%%%%%%%%%%%%%%%%%%封面与目录%%%%%%%%%%%%%%%%%%%%%%%%%%%%%%
\begin{titlepage}
  \begin{center}
  % Upper part of the page
  \includegraphics[width=0.25\textwidth]{imgs/logo.jpg}\\[1cm]
  \textsc{\LARGE Department of Automation}\\[1.5cm]
  \fangsong{\Large 示例模板}\\[0.5cm]
  % Title
  \hrulefill
  \\[0.8cm]{\centering \LARGE \heiti 文章标题}\\[0.4cm]
  \hrulefill
  \\[4cm]

  % Author and supervisor
  \begin{tabbing}       %tabbing  列表
    \hspace*{4.6cm} \= \hspace{2.6cm} \= \kill
    % \=     in tabbing environment, sets a tab stop
    % \kill  in a\tabbing environment, deletes previous line so tabs can be set without outputting text.
    % \>     in tabbing environment is a forward tab.

    \>{\fangsong\sihao\textbf {班\hspace{1cm}级 \ \ ：}}\>  {\centering\fangsong\sihao\textbf{~~~~~~~~自~~3~2}} \\
    \\
    \>{\fangsong\sihao\textbf {姓\hspace{1cm}名 \ \ ：}}\>  {\centering\fangsong\sihao\textbf{~~~~~~~~陈~昊~楠}}\\
    \\
    \>{\fangsong\sihao\textbf {学\hspace{1cm}号 \ \ ：}}\>  {\centering\fangsong\sihao\textbf{~~~~~~2013011449}}\\
    \\
    \>{\fangsong\sihao\textbf {指导教师 \ \ ：}}\>  {\centering\fangsong\sihao\textbf{~~~~~~~~~张~~~毅}} \\
  \end{tabbing}
  \vfill
  {\large \today}
  \end{center}
\end{titlepage}

\tableofcontents
\clearpage

%%%%%%%%%%%%%%%%%%%%%%%%%%正文部分%%%%%%%%%%%%%%%%%%%%%%%%%%%%%%%%%%

\begin{table}[htbp]
\centering
\pgfplotstabletypeset[
	every head row/.style={before row=\hline, after row=\hline},
	every last row/.style={after row=\hline},
	every first column/.style={column type/.add={|}{|}},
	every last column/.style={column type/.add={}{|}},
	columns/gforce/.style={column name=所加砝码质量(g)},
	columns/analog/.style={column name=输入通道示数},
]{data/weight.txt}
\caption{重力传感器所加砝码质量与输入通道示数}
\label{tab:weight}
\end{table}

\begin{figure}[htbp]
\centering
\begin{tikzpicture}
\begin{axis}[
    xlabel={砝码质量(g)},
    ylabel={模拟量输出},
    % legend pos=south east,
    % legend entries={Concrete,Linoleum},
    ]
  \addplot table [x=gforce,y=analog] {data/weight.txt};
\end{axis}
\end{tikzpicture}
\caption{重力传感器输出特性曲线}
\label{fig:weight}
\end{figure}

\end{document}
